
% ────────────────────────────────────────────────────────
\section{O Teorema Fundamental do Cálculo Integral}

% ────────────────────────────────────────────────────────
\subsection{Enquadramento Teórico}

\begin{thmbox}[title={Teorema Fundamental do Cálculo Integral}]
  Seja $f:[a,b]\to\mathbb{R}$ contínua.
  \begin{enumerate}
    \item \textbf{(Parte 1)} A função $G(x)=\displaystyle\int_a^x f(t)\,\dd t$
          é diferenciável em $(a,b)$ e $G'(x)=f(x)$.
    \item \textbf{(Parte 2)} Se $F$ for qualquer primitiva de $f$, então
          \[
            \int_a^b f(x)\,\dd x = F(b)-F(a) =: \Big[F(x)\Big]_a^b.
          \]
  \end{enumerate}
\end{thmbox}

% ────────────────────────────────────────────────────────
\subsection{Exemplos Resolvidos}

\begin{exbox}{5 --- Aplicação do Teorema Fundamental}
  Calcule $\displaystyle\int_0^{\pi} \sin x\,\dd x$.
\end{exbox}
\begin{solbox}
  Uma primitiva de $\sin x$ é $-\cos x$.
  \[
    \int_0^{\pi}\sin x\,\dd x = \Big[-\cos x\Big]_0^{\pi}
      = (-\cos\pi)-(-\cos 0) = 1+1 = 2.
  \]
\end{solbox}
